% Generated draft: requires mathematical and source review.
\documentclass[11pt]{article}
\usepackage[margin=25mm]{geometry}
\usepackage{fontspec}
\setmainfont{Latin Modern Roman}
\usepackage{amsmath,amssymb,mathtools,unicode-math}
\setmathfont{Latin Modern Math}
\usepackage{enumitem,needspace,xurl,hyperref}
\setlength{\parindent}{0pt}
\setlength{\parskip}{5pt}
\setlength{\emergencystretch}{3em}
\newcommand{\N}{\mathbb{N}}
\newcommand{\Z}{\mathbb{Z}}
\newcommand{\Q}{\mathbb{Q}}
\newcommand{\R}{\mathbb{R}}
\newcommand{\C}{\mathbb{C}}
\newcommand{\F}{\mathbb{F}}
\newcommand{\A}{\mathbb{A}}
\newcommand{\PP}{\mathbb{P}}
\newcommand{\E}{\mathbb{E}}
\newcommand{\eps}{\varepsilon}
\newcommand{\dd}{\,\mathrm{d}}
\begin{document}
\section*{UnsolvedMath ID: 30006652}
\textbf{Source problem code:} catalogue entry
\subsection{Definitions and mathematical statement}
Let \(\mathbb S\) denote the sphere spectrum. Following the supplied flatness condition, a spectrum \(A\) is flat over \(\mathbb S\) if its abelian group \(R=\pi_0(A)\) is flat over \(\Z=\pi_0(\mathbb S)\), and the maps induced by the sphere action are isomorphisms
\[
R\otimes_{\Z}\pi_n(\mathbb S)\xrightarrow{\ \cong\ }\pi_n(A)
\qquad(n\in\Z).
\]
Here flatness of \(R\) means that \(R\otimes_{\Z}-\) preserves exact sequences of abelian groups.

Equip \(R\) with a commutative ring structure extending its given abelian group structure. For each prime \(p\), write
\[
\widehat R_p=\varprojlim_{k\geq1}R/p^kR,
\]
and prescribe a Frobenius lift \(\varphi_p:R\to\widehat R_p\). Under the usual ring-theoretic interpretation, this is a ring homomorphism satisfying
\[
\rho_p(\varphi_p(r))=r^p+pR
\qquad(r\in R),
\]
where \(\rho_p:\widehat R_p\to R/pR\) is the canonical projection. The supplied excerpt does not explicitly define completion or Frobenius lift; these formulas state the conventional interpretation used in this draft. The unit convention and any additional compatibility conditions between different primes require confirmation from the full source.

An \(E_\infty\) ring structure means a multiplication and unit on a spectrum with the coherent commutativity and associativity encoded by an \(E_\infty\) operad. The supplied Proposition 7 associates to a flat \(E_\infty\) ring spectrum a canonical map
\[
A\longrightarrow (A^{\otimes p})^{tC_p}
\longrightarrow A^{tC_p}\simeq A_p^{\wedge}.
\]
Here \(\otimes\) is the smash product over \(\mathbb S\), \(C_p\) acts on \(A^{\otimes p}\) by cyclic permutation, and \((-)^{tC_p}\) denotes the Tate construction. In the displayed target \(A^{tC_p}\), the action on \(A\) is trivial. The source describes the resulting maps on \(\pi_0\) as the Frobenius lifts above; its completion identification is retained as a supplied assertion.

\textbf{Problem.} For every flat spectrum \(A\), characterize precisely the data
\[
\bigl(R,\text{ commutative ring structure},(\varphi_p)_{p\text{ prime}}\bigr)
\]
that are induced by an \(E_\infty\) ring structure on \(A\). More explicitly, determine necessary and sufficient conditions for there to exist an \(E_\infty\) ring spectrum \(B\) and an equivalence of underlying spectra \(e:B\simeq A\) such that \(\pi_0(e)\) transports both the ring structure and every induced Frobenius lift to the prescribed data on \(R\). For realizable data, determine how much of the \(E_\infty\) structure is determined by those data, including whether the extension is unique.

The uniqueness target remains unresolved at the level of formulation: the supplied record does not specify whether it asks for equivalence classes of unmarked ring spectra, equivalences compatible with the identification with \(A\), or a contractible space of extensions. These are distinct assertions and are not identified here.

As supplied background, Theorem 9 states that \(H_\infty\) structures on a flat spectrum are fully determined by the commutative ring on \(\pi_0\) and its Frobenius lifts. This is recorded as the source's partial result, without asserting that every prescribed family is realizable or that it determines an \(E_\infty\) structure. The original Question 8 asks broadly how much of the \(E_\infty\) structure is determined; the registry's reconstructed statement explicitly asks for an existence characterization and uniqueness.
\subsection{Short English statement}
For a spectrum flat over the sphere, determine which commutative ring structures on its zeroth homotopy group, together with a Frobenius lift at every prime, come from an \(E_\infty\) ring structure. Determine also the extent of uniqueness when such a structure exists. The source supplies an \(H_\infty\) partial result, while the precise meaning of uniqueness and some conventions require clarification.
\subsection{Sources}
\begin{itemize}
\item Supplied local UnsolvedMath registry record, permanent ID \(30006652\), title: \(E_\infty\) Structures from Frobenius Lifts; exact source problem code \texttt{catalogue entry}. The record attributes the question to an Oberwolfach workshop report, gives proposed year 2026, and names no individual proposer.
\item \emph{Cohomology of Finite Groups: Interactions and Applications} (2026), Oberwolfach Reports, pp.\ 541--610, DOI \texttt{10.4171/owr/2026/9}: \url{https://doi.org/10.4171/owr/2026/9}. The supplied excerpt identifies Question 8, Proposition 7, and Theorem 9. It attributes the observation underlying Proposition 7 to Nikolaus--Scholze.
\item Carmeli--Nikolaus--Yuan, \emph{Maps between spherical group rings}: \url{https://arxiv.org/abs/2405.06448}. Listed in the supplied literature sources; its specific bearing on the question is not detailed in the record.
\item Nikolaus--Scholze, \emph{On topological cyclic homology}: \url{https://doi.org/10.4310/ACTA.2018.v221.n2.a1}. Listed in the supplied literature sources and associated by the excerpt with the Frobenius construction.
\item Supplied literature assessment checked on \texttt{2026-08-22}: the report is said to determine \(H_\infty\) structures from the ring and Frobenius lifts while leaving extension and uniqueness of \(E_\infty\) structures unresolved. The registry labels the problem \texttt{partially\_solved}. This is supplied dated curation, not an independently verified assessment. The record also labels its reconstructed statement \texttt{corrected\_verified} and says it was checked against the official report; that verification claim is supplied metadata and has not been independently checked here.
\end{itemize}
\end{document}
